% ==STYLE-BEGIN==
\documentclass[UTF8,11pt]{ctexart}
\usepackage{amsmath,amssymb,amsthm}
\usepackage[a4paper,margin=2.4cm,top=2.3cm,bottom=2.3cm]{geometry}
\usepackage{fancyhdr}
\usepackage{enumitem}
\usepackage{hyperref}
\hypersetup{hidelinks}

% ---- 版式 ----
\linespread{1.15}
\setlength{\parskip}{0.35em}
\setlength{\parindent}{1.6em}
\setlist{nosep,leftmargin=2em}

% ---- 页眉 ----
\pagestyle{fancy}
\fancyhf{}
\fancyhead[L]{\small\coursename}
\fancyhead[R]{\small\thepage}
\renewcommand{\headrulewidth}{0.4pt}

% ---- 定理环境：中文用正体，标题加粗，按小节共用一条流水号 ----
\newtheoremstyle{cnthm}{0.7em}{0.7em}{\normalfont}{}{\bfseries}{}{0.5em}{}
\theoremstyle{cnthm}
\newtheorem{thm}{定理}[section]
\newtheorem{lem}[thm]{引理}
\newtheorem{prop}[thm]{命题}
\newtheorem{cor}[thm]{推论}
\newtheorem{defn}[thm]{定义}
\newtheorem{rem}[thm]{注}
\newtheorem{exa}[thm]{例}
\newtheorem{conj}[thm]{猜想}
\renewcommand{\proofname}{\bfseries 证明}
\providecommand{\coursename}{}
% ==STYLE-END==
\renewcommand{\coursename}{数学分析\,·\,1021}
\title{数学分析 1021}
\date{}
\begin{document}
\maketitle

\section{紧集}

\begin{defn}[紧集]
$S\subseteq\mathbb{R}$ 称为紧集，是指：若
\[
S\subseteq\bigcup_{i\in I}U_i,\qquad U_i\ \text{均为开集},
\]
则 $\exists i_1,\cdots,i_k$ s.t. $S\subseteq U_{i_1}\cup\cdots\cup U_{i_k}$。
\end{defn}

回顾有限覆盖定理：有界闭区间一定是紧集。

\begin{thm}[紧集的刻画]
$C$ 是 $\mathbb{R}$ 上的紧集 $\iff C$ 是有界的且 $C$ 是闭集。
\end{thm}

分析：已知有界闭区间一定是紧集，只需证明紧集一定是有界闭集。即设 $X\subseteq\mathbb{R}$ 是紧集，需要证明：1. $X$ 有界；2. $\mathbb{R}\setminus X$ 是开集。

\begin{proof}
（1）$X$ 有界。$\forall x\in X$，构造开区间 $(x-1,x+1)$，得到开覆盖
\[
X\subseteq\bigcup_{x\in X}(x-1,x+1).
\]
由紧集的定义，一定 $\exists x_1,\cdots,x_k\in X$，s.t.
\[
X\subseteq(x_1-1,x_1+1)\cup\cdots\cup(x_k-1,x_k+1),
\]
于是 $X$ 有界。
\end{proof}

\begin{proof}
（2）$\mathbb{R}\setminus X$ 是开集。若 $X=\varnothing$，则 $\mathbb{R}\setminus X=\mathbb{R}$ 是开集，结论平凡；以下设 $X\neq\varnothing$。取定任意 $y_0\in\mathbb{R}\setminus X$，则对每个 $x\in X$ 有 $x\neq y_0$，取 $\delta_x=\eta_x=\frac{|x-y_0|}{2}>0$，一定有
\[
\forall x\in X,\ \exists\delta_x,\eta_x\ \text{s.t.}\
(y_0-\delta_x,y_0+\delta_x)\cap(x-\eta_x,x+\eta_x)=\varnothing.
\]
这些邻域构成 $X$ 的一个开覆盖：
\[
X\subseteq\bigcup_{x\in X}(x-\eta_x,x+\eta_x).
\]
而 $X$ 是紧集，一定有有限的子覆盖，即
\[
\exists x_1,\cdots,x_k\ \text{s.t.}\
X\subseteq(x_1-\eta_{x_1},x_1+\eta_{x_1})\cup\cdots\cup(x_k-\eta_{x_k},x_k+\eta_{x_k}),
\]
且有
\[
(y_0-\delta_{x_j},y_0+\delta_{x_j})\cap(x_j-\eta_{x_j},x_j+\eta_{x_j})=\varnothing
\quad(j=1,\cdots,k).
\]
取 $\delta=\min\{\delta_{x_1},\cdots,\delta_{x_k}\}>0$（这里用到 $k\ge 1$，即 $X\neq\varnothing$）。由于 $\delta\le\delta_{x_j}$，有
\[
(y_0-\delta,y_0+\delta)\subseteq(y_0-\delta_{x_j},y_0+\delta_{x_j}),
\]
故 $(y_0-\delta,y_0+\delta)$ 与每个 $(x_j-\eta_{x_j},x_j+\eta_{x_j})$ 都不相交，即有
\[
(y_0-\delta,y_0+\delta)\cap X=\varnothing.
\]
这说明 $\mathbb{R}\setminus X$ 是开集，故 $X$ 是闭集。
\end{proof}

\section{极值定理}

\begin{thm}[极值定理]
设 $f:[a,b]\to\mathbb{R}$ 连续，则 $\operatorname{Im}f=f([a,b])$ 是有界闭区间；特别地，存在 $\xi,\eta\in[a,b]$ 使得
\[
f(\xi)=\inf_{[a,b]}f,\qquad f(\eta)=\sup_{[a,b]}f.
\]
\end{thm}

\begin{proof}
已知 $[a,b]$ 是紧集。「是区间」由介值定理（见 1019）保证；「有界且闭」这里不去用「有界闭 $\Rightarrow$ 紧」那一半，而是把 $[a,b]$ 的紧性直接推给 $\operatorname{Im}f$。

取 $\operatorname{Im}f$ 的任一开覆盖 $\{U_i\}_{i\in I}$，即
\[
\operatorname{Im}f\subseteq\bigcup_{i\in I}U_i,
\]
目标是找到 $i_1,i_2,\cdots,i_k$，s.t.
\[
\operatorname{Im}f\subseteq U_{i_1}\cup\cdots\cup U_{i_k}.
\]
由 $\operatorname{Im}f\subseteq\bigcup\limits_{i\in I}U_i$ 且
\[
f^{-1}(U_i)=f^{-1}(U_i\cap\operatorname{Im}f).
\]
由于 $f$ 连续，开集的原像仍为开集，所以 $f^{-1}(U_i)$ 是 $[a,b]$ 中的（相对）开集：严格地说 $f^{-1}(U_i)=O_i\cap[a,b]$，其中 $O_i$ 是 $\mathbb{R}$ 中的开集。取原像有
\[
[a,b]\subseteq\bigcup_{i\in I}f^{-1}(U_i).
\]
由有限覆盖定理，一定有
\[
\exists j_1,\cdots,j_k\ \text{s.t.}\
[a,b]\subseteq f^{-1}(U_{j_1})\cup\cdots\cup f^{-1}(U_{j_k}).
\]
全部拉回，有
\[
f([a,b])\subseteq U_{j_1}\cup\cdots\cup U_{j_k}.
\]
因此 $f([a,b])$ 是紧集，从而是有界闭区间，故 $\sup\operatorname{Im}f$、$\inf\operatorname{Im}f$ 都属于 $\operatorname{Im}f$。
\end{proof}

\section{介值定理与道路连通}

看完了极值定理的拓扑证明，不禁要问：上一节用分析的手段证过介值定理，那么对介值定理，是否也能用类似的手段简单地证出来？为此引入道路连通。

\begin{defn}[道路]
取 $S\subseteq\mathbb{R}$，$S$ 上的一条道路是指
\[
\alpha:[0,1]\to\mathbb{R}\ \text{连续},\qquad \alpha([0,1])\subseteq S.
\]
\end{defn}

\begin{defn}[道路连通]
如果 $\forall x,y\in S$，存在 $S$ 中的道路 $\alpha:[0,1]\to\mathbb{R}$ 使得 $\alpha(0)=x$、$\alpha(1)=y$，则称 $S$ 为（道路）连通的。
\end{defn}

\begin{conj}
$\mathbb{R}$ 中道路连通的子集必为区间。
\end{conj}

在上述猜想成立的情况下，由介值定理可知，连续函数会把一个区间映成一个区间，猜想保证了区间必然道路连通。所以对连续函数，猜想它会把一个道路连通的集合映成一个道路连通的集合，即

\begin{thm}[连续像保持道路连通]
对于 $f:I\to\mathbb{R}$，若 $I$ 道路连通且 $f$ 连续，那么 $\operatorname{Im}f$ 也是道路连通的。
\end{thm}

\begin{proof}
即 $\forall y_1,y_2\in\operatorname{Im}f$，找一个连续映射 $\beta:[0,1]\to\mathbb{R}$，使得
\[
\beta(0)=y_1,\quad \beta(1)=y_2,\quad \beta([0,1])\subseteq\operatorname{Im}f.
\]
取 $x_1,x_2\in I$ 使 $f(x_1)=y_1$，$f(x_2)=y_2$。对于 $I$，因为其道路连通，一定有其中的道路 $\alpha:[0,1]\to\mathbb{R}$ 使得 $\alpha(0)=x_1$，$\alpha(1)=x_2$。$f$ 连续，连续函数的复合也是连续的，因此 $f\circ\alpha$ 就是所要的 $\beta$。
\end{proof}

\end{document}
